% English translation, made from our own transcription (original.tex), not from any existing
% translation. Now conformed to the 1694 Acta original (see original.tex header). Every math
% expression, \origpage marker (336-338) and the \rmfigure image path match original.tex; only the
% prose is translated. Notes for the reviewer:
%  * This is the September 1694 sequel that first names the lemniscate. The author's embedded French
%    phrase "d'un nœud de ruban" is kept untranslated (his own gloss), with "Gallis" rendered "as the
%    French call it".
%  * The lettered footnotes (a)-(f) of the 1744 Opera (Cramer's apparatus) are absent from the Acta
%    and not present here; the Opera's closing paragraph "Caeterum quantitas ..." is likewise an
%    Opera addition absent from the Acta and not translated.
%  * The author's "&c." is given as "etc." and the Latin "&" (et) as "and"; his own asides use round
%    parentheses, translator interpolations use square brackets (R9); math (incl. the ":" division
%    sign) is copied verbatim from original.tex.
\documentclass{article}
\usepackage{readmasters}
\begin{document}

\origpage{336}
\section*{Jac.\ B. Construction of the Curve of uniform Approach \& Recession}
\subsection*{by means of the rectification of a certain algebraic Curve: an Addendum to the recent solution of the Month of June}

There are chiefly three modes of constructing mechanical, or transcendental, curves. The first, but
little suited to practice, is done through the quadratures of curvilinear spaces. Better is that which
is set up through the rectifications of algebraic curves; for curves can be rectified in practice more
accurately and expeditiously, by means of a thread or a little chain wrapped around them, than spaces
can be squared. In the same rank I hold those constructions which are carried out without any
rectification and quadrature, by the mere description of some mechanical curve whose points, though
not all, yet infinitely many and as close together as one pleases, can be found geometrically, such
as the Logarithmic curve is wont to be, and any others there may be of that kind. But the best mode,
wherever it can be had, is that which is carried out by means of some curve which nature herself,
without art, produces by a certain most swift and, as it were, instantaneous motion at the beck of
the Geometer; whereas the preceding modes require curves whose

\origpage{337}
delineation, whether by continuous motion or by the finding of several points, being set up by the
Craftsman, commonly turns out either too slow or too much impeded: thus those constructions of
Problems which suppose the quadrature of the Hyperbola or the description of the Logarithmic curve I
judge, other things being equal, to be inferior to those which are carried out by means of the
Catenary, that is, of the curve which a suspended chain assumes of its own accord sooner than you
could set the first hand to describing the others.

\emph{The Curve of uniform Approach and Recession}, for the sole recommendation of the Gentleman who
proposed it, seems to deserve that we should take pains to furnish it with constructions according to
all three modes. The construction of the first mode, by quadratures, which I passed over in the recent
solution for the reason already stated, can be set up by taking, in \emph{Fig.}~3,\ \emph{Tab.}~VI,
the straight line $EN$ equal, not, as there, to $Au$ itself, but to $iu$; and again the rectangle
$AO$ made equal to the curvilinear area $AEN$; for its side $AP$ will be adequate to the recent curve
$AQ$, so that, taking $A\alpha$ as the third proportional to $AB$ and the now-found $AP$, the point
$\alpha$ is had on the desired Curve. The construction of the third mode, which would be done by means
of the Elastic line $AQ$, which alone I touched upon there, would without doubt be the best, if nature
had anywhere made tensions simply proportional to the stretching forces: but since perhaps no Spring
is found which observes this law of tension precisely; nor, if one were found, would that be certain;
therefore it is not safe enough to trust to it either; and it is preferable to have recourse to the
second mode of construction, and to seek an algebraic Curve by whose rectification we may attain our
aim. Such a one, as if by wish, presents itself: a Curve of four dimensions which is expressed by this
equation $xx + yy = a\sqrt{xx - yy}$, and which, set about the axis $BG\ (2a)$, refers in form to a
reclining figure of the numeral eight $\infty$, or to a band folded into a knot, that is, into a
lemniscus, \emph{d'un nœud de ruban} as the French call it. For of this curve both the one half is
[equal to] the Elastic curve $AR\varphi$, and, above the node $A$, over the indeterminate interval
$AE$ cut off, the smaller portion of the half [equals] the smaller $AQ$, the larger [equals] the
larger Elastic portion $\varphi RQ$. Whence too the reason is established for employing it in the
construction of the proposed Problem.

\rmfigure{figures/fig-3.png}{Fig.~3.}{Fig. 3 (Tab. VI): the elastic curve $AR\varphi$, the lemniscate
looped through the node $A$, the isochronal curve of uniform approach and recession, and the
auxiliary circle, quadrants and construction lines.}

For the rest, just as infinitely many mechanical lines depend on the rectification of the circular or
parabolic curve, or on the logarithmic description; so too, besides the curve of approach and
recession of infinitely many others, the construction of the Elastic curve itself (as will appear

\origpage{338}
below) is in part derived from the rectification of the said algebraic curve of four dimensions. And
I would dare to assert that, in the matter of the construction of mechanical curves, this one
immediately follows the preceding among the rest; so much so that a construction which succeeds by
none of the earlier ones is to be attempted next either by the rectification of this Curve, or of a
Hyperbolic or an Elliptic line, or of two at once. The reason for this is that, after the differential
formulas $aa\,dz : \sqrt{aa - zz}$, $zz\,dz : \sqrt{aa - zz}$, $aa\,dz : \sqrt{zz - aa}$,
$zz\,dz : \sqrt{zz - aa}$, which are integrated by means of the quadrature of the Circle and the
Hyperbola, the simplest by far seem to be these expressions, $zz\,dz : \sqrt{z^{4}-a^{4}}$,
$aa\,dz : \sqrt{z^{4}-a^{4}}$, $aa\,dz : \sqrt{a^{4}-z^{4}}$, $zz\,dz : \sqrt{a^{4}-z^{4}}$, and the
like; of which the first is integrated by means of a Hyperbolic line, the second and third by the
rectification of our lemniscate curve, the fourth by that of the same and of an Elliptic curve. For
indeed if the indeterminate $z$ be applied to an equilateral Hyperbola (whose transverse axis is
$2a$) from its centre, the arc intercepted between the vertex and the applied line will satisfy the
first formula. And if the same $z$ (or the third proportional to $z$ and $a$, where $z$ is greater
than $a$) be subtended from the node of the lemniscate Curve, the subtended arc will serve for the
two middle ones. And if, finally, from the centre of an Ellipse (whose minor axis is $2a$, the major
$2a\sqrt{2}$) the same $z$ be cut off on the minor axis and through the point of section a straight
line be drawn parallel to the major, the elliptic portion intercepted between the parallels will
designate it, the corresponding portion of the lemniscate Curve being cut away, the integral of the
fourth formula $zz\,dz : \sqrt{a^{4}-z^{4}}$. Whence, since the same line in the cited figure,
$AE = z$ being given, denotes the element of the applied line $AP$ or $EQ$ in the Elastic curve, it is
plain that this applied line equals the difference of two portions of the Elliptic and the Lemniscate
curves, and therefore it is manifest how the Elastic curve can be constructed by the rectification of
both. And since the portion of the Lemniscate, as I noted above, is adequate to the Elastic portion
$AQ$, it is clear also that the Elliptic curve itself is adequate to the aggregate $AQ + AP$; which
stands as a not inelegant property of these Curves.

\end{document}
